\chapter{Modelos matemáticos de sistemas dinámicos lineales}

\section{Generalidades}

Existen diversas formas de modelar un sistema dinámico en el rango lineal.

\begin{figure}[H]
    \centering
    \includegraphics[scale=0.7]{modelos.png}
    \caption{Relaciones entre distintas representaciones de modelos de sistemas dinámicos.}
\end{figure}

\section{Principios fundamentales de la mecánica estructural}

Dominio $\Omega$ tiene una frontera $\partial \Omega = \Gamma_u \cup \Gamma_t$, donde $\Gamma_u$ es la frontera con desplazamientos prescritos, y $\Gamma_t$ es la frontera con tracciones prescritas.

\begin{figure}[H]
    \centering
    \includegraphics[width=4in]{dominio.png}
    \caption{Problema de valores de contorno en mecánica estructural.}
\end{figure}

\begin{enumerate}
    \item Cinemática y compatibilidad
    \item Ecuaciones de balance
    \item Leyes constitutivas de los materiales
    \item Geometría y propiedades estructurales
    \item Condiciones de borde y condiciones iniciales
\end{enumerate}

\subsection{Forma fuerte}

La ecuación de equilibrio en forma fuerte, que describe el equilibrio dinámico en un continuo elástico, donde:

\begin{align}
    \sigma_{ij,j} + b_i = \rho\frac{\partial^2 u_i}{\partial t^2}, & \qquad \forall \Omega\\
    u_i = q_i, & \qquad \forall \Gamma_u\\
    \sigma_{ij} n_j = t_i, & \qquad \forall \Gamma_t\\
    u_i(x,0) = u_{0i}(x), & \qquad \forall \Omega\\
    \frac{\partial u_i}{\partial t}(x,0) =\frac{\partial u_{0i}}{\partial t}(x), & \qquad \forall \Omega\\
\end{align}

donde: $\sigma_{ij,j}$ es la divergencia del tensor de esfuerzos ($\sigma_{ij}$); $b_i$ son las fuerzas volumétricas; $\rho \frac{\partial^2 u_i}{\partial t^2}$ es el término inercial, con $\rho$ la densidad y $u_i$ el desplazamiento.

\subsection{Forma débil}

Lagrangiano: $\mathcal{L} = T-V$

Acción: $S = \int_{t_1}^{t_2} \mathcal{L}dt$

Principio de mínima acción: 

\begin{equation}
    \delta S = 0
\end{equation}

Utilizando calculo variacional se obtiene la forma débil del problema de elastodinámica, utilizando notación directa (Hughes, 2000):

\begin{equation}
    (w,\rho \ddot{u}) + a(w,u) = (w,b) + (w,t)_\Gamma
\end{equation}

\subsection{Discretizacion}

\begin{itemize}
    \item Rayleigh-Ritz
    \item Modos asumidos
    \item Métodos residuales (por ejemplo, Galerkin)
\end{itemize}

\section{Ecuaciones de movimiento}

Forma matricial obtenida después de la etapa de semidiscretización espacial (tiempo contínuo). El sistema dinámico se escribe finalmente como un sistema de ecuaciones diferenciales ordinarias (ODEs).

\begin{equation}
    M \ddot{u} + C \dot{u} + K u = F(t)
\end{equation}

donde $u = \{u_1,u_2, \ldots,u_n\}^T$ son los grados de libertad del sistema.

\section{Funciones de respuesta en frecuencia (FRF)}

\emph{En desarrollo}

\section{Función de transferencia (TF)}

\subsection{Sistema 1GDL}

\begin{figure}[H]
    \centering
    \includegraphics[scale=0.7]{edm.png}
\end{figure}

\underline{Problema}: Dada la EDM de un sistema estructural 1GDL, determinar su función de transferencia $H(s)$\\ 

\noindent Notación: 

\begin{align*}
    U(s) &= \mathcal{L}\{u(t)\}~: \text{respuesta (dominio complejo)}\\
    F(s) &= \mathcal{L}\{f(t)\} ~: \text{input (dominio complejo)}
\end{align*}

Asumiendo reposo: $u(0)=\dot{u}(0)=0$

\begin{align*}
    m\ddot{u}+c\dot{u}+ku &= f(t) ~/\color{red}{\mathcal{L}\{\cdot\}}\\
    m\mathcal{L}\{\ddot{u}\}+c\mathcal{L}\{\dot{u}\}+k\mathcal{L}\{u\} &= \mathcal{L}\{f\}\\
    ms^2U(s)+csU(s)+kU(s) &= F(s)\\
    (ms^2+cs+k)U(s) &= F(s)
\end{align*}
\begin{align*}
    \Longrightarrow H(s) = \frac{U(s)}{F(s)}&=\frac{1}{ms^2+cs+k} ~\text{: TF del sistema de 1GDL}
\end{align*}

Aparte: Dada la entrada $F(s)$ y la TF $H(s)$, la salida es $U(s) = H(s)F(s)$
\begin{figure}[H]
\centering
\includegraphics[scale=0.7]{hs.png}
\end{figure}

\subsection{Sistema MDGL}

Ahora considere la EDM de un sistema MGDL:

\begin{equation}
    M \ddot{u}(t )+ C \dot{u}(t) + K u(t) = F(t)
\end{equation}

Donde: $u(t) \in \mathbb{R}^{n \times 1}$, es el vector de grados de libertad; $M \in \mathbb{R}^{n \times n}$, $C \in \mathbb{R}^{n \times n}$, y $K \in \mathbb{R}^{n \times n}$ son las matrices de masas, amortiguamiento y rigidez, respectivamente; $F(t) \in \mathbb{R}^{n \times 1}$ es el vector de fuerzas externas.\\

\noindent Notación: 

\begin{align*}
    \mathcal{L}\{u(t)\} &= U(s)\\
    \mathcal{L}\{\dot{u}(t)\} &= s U(s) - u(0)\\
    \mathcal{L}\{\ddot{u}(t)\} &= s^2 U(s) - s u(0) - \dot{u}(0)\\
    \mathcal{L}\{f(t)\} &= F(s)
\end{align*}

Entonces, la función de transferencia (\emph{transfer function, TF}) del sistema dinámico se obtiene aplicando la transformada de Laplace sobre la EDM, asumiendo condiciones iniciales de reposo ($u(0)= 0 \wedge \dot{u}(0)=0$):

\begin{align*}
    M \ddot{u}(t) + C \dot{u}(t) + K u(t) &= F(t) ~/\mathcal{L}\{\cdot\}\\
    %\Longrightarrow s^2 M U(s) + s C U(s) + K U(s) &= F(s)\\
    \Longrightarrow \left[s^2M+sC+K \right]U(s)&= F(s) \\
     \Longrightarrow U(s)&= \underbrace{\left[s^2M+sC+K \right]^{-1}}_{H(s)} F(s)
\end{align*}


\begin{align*}
    \tend{F}(s) &= \mathcal{L}\{\tend{f}(t)\} ~: \text{input en dominio complejo (excitación)}\\
    \tend{U}(s) &= \mathcal{L}\{\tend{u}(t)\}~: \text{output en dominio complejo (respuesta)}\\
    \tenq{H}(s) &= \left[s^2M + s\tenq{C}+\tenq{K} \right]^{-1} ~:\text{TF para un sistema MGDL}
\end{align*}

Notar que para encontrar los polos se debe resolver las raíces polinomio:
\begin{align*}
    s^2\tenq{M}+s\tenq{C}+\tenq{K} = \tend{0} ~/ \text{$\cdot\tend{\phi}$ yasumiendo $\tenq{C}=\tenq{0}$} \\
    \Longrightarrow (s^2\tenq{M}+\tenq{K})\tend{\phi} = \tend{0} ~\text{: Problema de valores propios generalizado}
\end{align*}


\section{Representación en variables de estado}

La representación espacio-estado reduce la EDM a ecuaciones diferenciales ordinarias cuyas matrices contienen toda la información dinámica del sistema.

\begin{align*}
    \text{EDM: }~\tenq{M}\tend{\ddot{u}}+\tenq{C}\tend{\dot{u}}+\tenq{K}\tend{u}&=\tenq{G}\tend{p}(t) 
\end{align*}

\begin{align*}
\left .
\begin{array}{rr}
\tend{\dot{x}}=\tenq{A_s}\tend{x}+\tenq{B_s}\tend{p}\\
\tend{y}=\tenq{C_s}\tend{x}+\tenq{D_s}\tend{p} 
\end{array}
\right \} \Longrightarrow 
\begin{array} {|r|r|}
  \hline A & B \\
  \hline C & D \\
  \hline 
\end{array}
\end{align*}

\begin{figure}[H]
\centering
\includegraphics[scale=0.7]{SS.png}
\end{figure}

\underline{Ecuación de estado}

\begin{align*}
    \tend{\dot{x}}=\tenq{A_s}\tend{x}+\tenq{B_s}\tend{p}
\end{align*}

Donde:
\begin{itemize}
    \item Vector de Estado: $\tend{x}=\begin{bmatrix}\tend{u} \\ \tend{\dot{u}}\end{bmatrix} \Longrightarrow \tend{\dot{x}}=\begin{bmatrix}\tend{\dot{u}}\\ \tend{\ddot{u}}\end{bmatrix}$
    \item Matriz de estado: $\tenq{A_s}=\begin{bmatrix}\tenq{0} & \tenq{I}\\ -\tenq{M}^{-1}\tenq{K} & -\tenq{M}^{-1}\tenq{C}\end{bmatrix}$
    \item Matriz de influencia de excitación: $\tenq{B_s}=\begin{bmatrix}\tenq{0}\\ \tenq{M}^{-1}\tenq{G}\end{bmatrix}$
\end{itemize}

\underline{Ecuación de respuesta} \\

\begin{align*}
    \tend{y}=\tenq{C_s}\tend{x}+\tenq{D_s}\tend{p}
\end{align*}

Donde (este es un ejemplo):
\begin{itemize}
    \item Vector de respuesta: $\tend{y}=\begin{bmatrix}\tend{u} \\ \tend{\dot{u}}\\ \tend{\ddot{u}}\end{bmatrix}$

    \item $\tenq{C_s}=\begin{bmatrix}\tenq{I} & \tenq{0}\\
    \tenq{0} & \tenq{I}\\ 
    -\tenq{M}^{-1}\tenq{K} & -\tenq{M}^{-1}\tenq{C}\end{bmatrix}$
    \item ``Feedthrough'': $\tenq{D_s}=\begin{bmatrix}\tenq{0}\\
    \tenq{0}\\ \tenq{M}^{-1}\tenq{G}\end{bmatrix}$
\end{itemize}


\subsection{Relación entre State-Space y TF}

Asumiendo condiciones iniciales de reposo: $\tend{x}(0) = \tend{0} \longrightarrow \text{Estados}$\\

Entonces, aplicando la transformada de laplace a la ecuación de estado:
\begin{align*} \label{eq:ss2tf_1}
    \tend{\dot{x}}&=\tenq{A_s}\tend{x}+\tenq{B_s}\tend{p} ~/\mathcal{L}\{\cdot\}\\
    \Longrightarrow s\tend{X}(s)&= \tenq{A_s}\tend{X}(s)+\tenq{B_s}\tend{P}(s)\\
    \Longrightarrow (s\tenq{I}-\tenq{A_s})\tend{X}(s)&=\tenq{B_s}\tend{P}(s)
\end{align*}
\begin{align}
    \Longrightarrow \tend{X}(s)&=(s\tenq{I}-\tenq{A_s})^{-1}\tenq{B_s}\tend{P}(s)
\end{align}

Luego, aplicando la transformada de Laplace a la ecuación de respuesta:

\begin{align*}
    \tend{y}(t)=\tenq{C_s}\tend{x}+\tenq{D_s}\tend{p} ~/\mathcal{L}\{\cdot\}\\
\end{align*}
\begin{align}\label{eq:ss2tf_2}
    \Longrightarrow \tend{Y}(s)&=\tenq{C_s}\tend{X}(s)+\tenq{D_s}\tend{P}(s)
\end{align} 

Reemplazando la ecuación \ref{eq:ss2tf_1} en \ref{eq:ss2tf_2}:
\begin{align*}
    \tend{Y}(s)&=\tenq{C_s}(s\tenq{I}-\tenq{A_s})^{-1}\tenq{B_s}\tend{P}(s)+\tenq{D_s}\tend{P}(s)\\
    \tend{Y}(s)&=[\tenq{C_s}(s\tenq{I}-\tenq{A_s})^{-1}\tenq{B_s}+\tenq{D_s}]\tend{P}(s)\\
    \Longrightarrow \tenq{H}(s)&=\tenq{C_s}(s\tenq{I}-\tenq{A_s})^{-1}\tenq{B_s}+\tenq{D_s}
\end{align*}

\underline{Aparte}: Asumiendo $\tenq{D_s}=0$:
\begin{align*}
    \tenq{H}(s)&=\tenq{C_s}(s\tenq{I}-\tenq{A_s})^{-1}\tenq{B_s}\\
    \tenq{H}(s)&=\frac{1}{\det(s\tenq{I}-\tenq{A_s})}\tenq{C_s}\text{ adj}(s\tenq{I}-\tenq{A_s})\tenq{B_s} \\
    \Longrightarrow\det(s\tenq{I}-\tenq{A_s})&=0 ~\text{: Problema de valores propios de }\tenq{A_s}
\end{align*}

Por lo tanto, los Polos se pueden obtener por las raíces del polinomio del denominador de $\tenq{H}(s)$, o, por los valores propios de $\tenq{A_s}$.




% \section{Implementación computacional}


